\documentclass[11pt]{article}
\usepackage[T1]{fontenc}
\usepackage{lmodern}
\usepackage[margin=1in]{geometry}
\usepackage{amsmath,amssymb,amsthm,mathtools}
\usepackage{booktabs,tabularx,array}
\usepackage{microtype}
\usepackage{xcolor}
\usepackage[colorlinks=true,allcolors=blue!45!black]{hyperref}
\usepackage{fancyhdr}
\hypersetup{pdftitle={A local moment refinement near the quartic zero-free boundary},pdfauthor={Zhongpai Gao},pdfsubject={Research draft: a proposed local moment and an exact conditional endpoint}}
\pagestyle{fancy}
\fancyhf{}
\fancyhead[L]{\small Local moment refinement}
\fancyhead[R]{\small Research draft \textbullet\ 11 October 2026}
\fancyfoot[C]{\thepage}
\setlength{\headheight}{14pt}
\setlength{\emergencystretch}{2em}
\setlength{\parskip}{3pt}
\allowdisplaybreaks[1]
\newtheorem{proposition}{Proposition}[section]
\newtheorem{lemma}[proposition]{Lemma}
\theoremstyle{definition}
\newtheorem{hypothesis}[proposition]{Hypothesis}
\newtheorem{remark}[proposition]{Remark}
\newcommand{\Q}{\mathbf Q}
\newcommand{\Z}{\mathbf Z}
\newcommand{\C}{\mathbf C}
\newcommand{\OO}{\mathcal O}
\newcommand{\eps}{\varepsilon}
\newcommand{\pos}[1]{(#1)_{+}}
\newcommand{\code}[1]{\texttt{\small #1}}
\title{A local moment refinement near the\protect\\quartic zero-free boundary\\[6pt]
\large A proposed analytic estimate and an exact conditional endpoint}
\author{Zhongpai Gao\thanks{Prepared with AI assistance in research, computation, formalization, and exposition. The upstream mathematical contributions are credited in Section~\ref{sec:credits}. This draft has not received independent mathematical certification.}}
\date{11 October 2026 \quad---\quad Draft 2}

\begin{document}
\maketitle
\begin{abstract}
We investigate a local modification of the moment estimate used in Collet's
quartic zero-free argument. The proposed estimate reduces the coefficient
of the total prime-slot length from $9/2$ to $2249/500=4.498$ in a small,
explicit range of the two ordinary summation lengths. The idea is to
treat two types of transformed rows differently: rows with smaller
buffered zero parameters use the existing moment after inexpensive slot
removal, whereas rows with larger parameters are counted jointly with
their common arithmetic labels. We explain the resulting exponent
calculations and give an exact, computer-assisted scalar certificate for
the candidate boundary $4374433/5000000=0.8748866$.
This number is strictly below the exact quartic comparison boundary.
The analytic local moment, its full application to the physical numerator,
and the all-family nonvanishing conclusion remain unproved in this draft.
Nine compiled Lean components establish intermediate statements under
their stated hypotheses; they do not establish that conclusion.
\end{abstract}

\section{Introduction and the status of the result}

A fixed zero-free half-plane is a statement that an $L$-function has no
zeros to the right of a specified vertical line. For the families under
discussion, lowering that line is an improvement. The eventual aim of
this project is a strict half-plane $\Re s>\theta$ for ordinary Dirichlet
$L$-functions, including the Riemann zeta function, and for finite-order
Hecke $L$-functions of trivial infinite type over $\Q(\sqrt{-3})$.
Principal poles are excluded. A bound near $0.875$ is a fixed-strip
improvement; it does not approach a proof of the Riemann hypothesis by
itself.

Our starting point is the pinned Cycle25 manuscript of Collet
\cite{collet}, which builds on OpenAI's $7/8$ framework \cite{openai}
and Levy's weighted numerator \cite{levy}. Its exact comparison boundary is
\begin{equation}\label{eq:quartic}
 b_0=\frac{11}{12}-\frac{\ell_0}{4}
     =0.8748866470399641\ldots,
 \qquad
 P(\ell_0)=0,
\end{equation}
where
\[
 P(X)=927X^4-3135X^3+2433X^2+275X-100,
 \qquad \ell_0\in[1/6,1/5].
\]
Throughout this paper, ``comparison boundary'' refers to this particular
pinned source, rather than to a claim about the current leaderboard.

The proposed modification is deliberately local. In the source proof,
two ordinary character sums and a product of short prime sums share a
moment budget. Write their logarithmic lengths as $n_1,n_2,z$ and the
row budget as $M$. The existing estimate is available with
$n_1+n_2+(9/2)z\le M$. We seek to replace $9/2$ by $4.498$ only when
the lengths lie near the configuration used by the endpoint argument.
Because $z>0$, a smaller coefficient permits a slightly larger prime
product at the same row budget. The weighted numerator can then gain a
small amount.

The proposed endpoint is
\begin{equation}\label{eq:theta}
 \theta=\frac{4374433}{5000000}=0.8748866.
\end{equation}
Its distance below \eqref{eq:quartic} is about $4.704\times10^{-8}$.
The size of this gain makes exact arithmetic essential. More
importantly, a favorable scalar inequality is useful only if it belongs
to a valid estimate for the original arithmetic sums.

There are three distinct levels of assertion in this draft. The
elementary inequalities and the exact endpoint certificate are established
within their stated algebraic scope. The intermediate Lean results have
been compiled against cached dependencies. The complete analytic estimate
and its nonvanishing application are proposed, with outstanding
obligations listed in Section~\ref{sec:gaps}. In particular, we do not
state \eqref{eq:theta} as a proved zero-free boundary.

For a first reading, Sections~\ref{sec:idea} and~\ref{sec:endpoint} explain
the mechanism and its numerical consequence. Sections~\ref{sec:high}
and~\ref{sec:interfaces} retain the arithmetic details that a source
review must check.

\section{The proposed local moment}\label{sec:moment}

\subsection{Objects and normalization}
We use the natural moment family of \cite[Appendix B3.4]{collet} over
$\OO=\Z[\omega]$, where $\omega=e^{2\pi i/3}$. Write $N$ for the
ideal norm, or the norm of an element. Set $U=Z^M$, with $Z>1$ and
$M>0$. At the parent stage the character acting on a summation variable
$n$ has the form
\[
 \psi_k(n)=\tau(n)\chi_n(k),
\]
where $\tau$ is a fixed finite ray character and $\chi_n(k)$ is the
source's sextic residue character, with its prescribed zero extension.
The ray modulus is fixed independently of $U$. No additional parent
character of positive moving width is allowed in our proposed estimate.
Such factors can occur in transformed children and must still be paid.

Let $S_{k,1}(Z^{n_1})$ and $S_{k,2}(Z^{n_2})$ denote the source's
centrally normalized ordinary character sums. The two smooth profiles,
real norm powers, and height parameters may be independent. A prime
slot has the form
\begin{equation}\label{eq:slot}
 Q_{k,i}=Z^{-z_i/2}
 \sum_{\mathfrak p\in\mathcal P_i^{\rm can}}
 \psi_k(\mathfrak p)\nu_i(\mathfrak p)
 W\!\left(\frac{N\mathfrak p}{Z^{z_i}}\right)
 \left(\frac{N\mathfrak p}{Z^{z_i}}\right)^{-\sigma_i+iv_i}.
\end{equation}
Here $\mathcal P_i^{\rm can}$ is the genuine canonical good-prime pool,
$\nu_i$ is an allowed relative ray character, and the slots share the
one smooth annular profile $W$. Put $z=\sum_i z_i$ and
$Q_{k,I}=\prod_{i\in I}Q_{k,i}$. Formula~\eqref{eq:slot} does not
permit arbitrary prime subsets or independently chosen slot profiles.

An additional common puncture is allowed: all factors may be restricted
to arguments coprime to a fixed nonzero mask of norm at most $U^B$.
It stays fixed during the current row sum but may depend on labels
already frozen in a preceding transformation. The underlying natural
zeros must also stay in place. A canceled character phase at a nonunit
does not turn the original value zero into one.

The word \emph{eligible} below means the literal source predicate
excluding all \code{FixedInducingRow} presentations. It cannot be
replaced by the simpler condition that one selected primitive character
is nonprincipal. The exceptional presentations have separate terms in
the full argument.

\begin{hypothesis}[Local physical moment, not yet proved]\label{hyp:local}
Use the preceding family, with all its natural masks and canonical slot
conditions. Suppose
\begin{equation}\label{eq:window}
 \frac{41M}{100}\le n_i\le\frac{11M}{25},\qquad
 \frac{3M}{100}\le z\le\frac M{25},\qquad
 n_1+n_2+Cz\le M,\qquad C=\frac{2249}{500}.
\end{equation}
For every $\eps>0$, the proposed estimate is
\begin{equation}\label{eq:moment}
 \sum_{k\ {\rm eligible}}
 \left|S_{k,1}(Z^{n_1})S_{k,2}(Z^{n_2})Q_{k,I}\right|^2
 \ll Z^{M+\eps}\mathcal H_J.
\end{equation}
The source's radial row weight is understood. The factor $\mathcal H_J$
denotes its finite smooth-seminorm and polynomial height allowance,
retaining independent heights. The assertion also includes the four
annular derivative-profile pairs and bounded scale-box recovery used by
the physical numerator. All quantifiers, including selection of a common
fine mesh before the later arithmetic data, are as in the source family.
\end{hypothesis}

The notation imports a substantial existing analytic framework. This
draft explains a modification of that framework; it is not a replacement
construction of the ordinary $L$-functions or of the source moment family.
The exact connection between \eqref{eq:moment} and the formal physical
state is one of the missing proofs.

\section{Why a local improvement might close}\label{sec:idea}

There are two potential difficulties with improving a moment estimate.
First, a transformed child need not remain in the same local length
window as its parent. Second, counting a character family after freezing
its arithmetic labels may be too expensive. The proposed argument treats
these difficulties separately.

Apply the two transformations in the existing moment proof and classify
the resulting nonexceptional rows by a buffered zero parameter $\delta$.
Roughly, the associated real zero location is $(1+\delta)/2$; the
buffer and finite character family are those of the source. When
$\delta\le7/10$, pointwise prime estimates make it cheap to remove
enough slots to recover the existing global moment condition. For
$\delta\ge7/10$, those pointwise estimates are more expensive, but the
rows are sparse. The proposed saving comes from counting each such row
\emph{together with} its common-support labels.

The split is made after the transformations. No Poisson formula is
applied to a zero-bin indicator. Also, the smaller-parameter branch calls
the existing moment rather than calling Hypothesis~\ref{hyp:local} on a
shortened child. Thus there is no claimed local induction whose length
window must somehow survive indefinitely.

\subsection{The size of the budget defect}
Define
\begin{equation}\label{eq:budget}
 A=n_1+n_2+z,\qquad
 g=\pos{n_1+n_2+\tfrac92z-M},\qquad
 a=\frac M{2000},\qquad \sigma=3a.
\end{equation}
The auxiliary amplifier has prime norms on the scale $Z^a$, disjoint
from the live slot windows after the mesh is chosen sufficiently fine.
Its length $a$ is fixed relative to $M$; it is not an arbitrarily small
final error.

\begin{lemma}[Elementary parent bounds]\label{lem:parent}
Under \eqref{eq:window},
\[
 0\le g\le\frac M{12500},\qquad
 \frac{85M}{100}\le A\le\frac{44753M}{50000}<\frac{9M}{10}.
\]
\end{lemma}
\begin{proof}
The difference of the two capacity coefficients is
$9/2-C=1/500$. Hence $g\le z/500\le M/12500$.
The lower bound for $A$ follows by adding the lower length bounds.
For the upper bound, use
\[
 A\le M-(C-1)z
 \le M-\frac{1749}{500}\frac{3M}{100}
 =\frac{44753M}{50000}.
\]
\end{proof}

\subsection{Rows with smaller zero parameters}
The intended child calculation gives a global-base defect at most
$g+6a+\eta$, where $\eta$ collects mesh, support, and buffering losses.
Deleting slots of total width $d_z$ lowers that defect by $(9/2)d_z$.
On a row whose buffered parameter is at most $k_0=7/10$, the proposed
pointwise deletion cost in squared norm is at most $k_0d_z$, plus the
reserved errors. The existing moment is then applied to the remaining
canonical state. If no slots remain, the existing zero-slot moment is
used instead.

The relevant scalar margin is elementary. Since the source child saving
is at least $a$, and $g\le M/12500<M/10000$,
\begin{equation}\label{eq:lowreserve}
 a-\frac{k_0}{9/2}(6a+g)
 \ge\frac M{2000}-\frac7{45}
       \left(\frac{3M}{1000}+\frac M{10000}\right)
 =\frac{4M}{225000}>0.
\end{equation}
One greedy prefix removes the necessary slots; its overshoot is one
mesh increment. The errors must be chosen smaller than the fixed margin
in \eqref{eq:lowreserve}. The compiled components support pointwise
removal and a restricted incumbent estimate, but applying them to all
actual transformed states is still required.

\section{Counting rows together with their labels}\label{sec:high}

\subsection{From sixth-power-free rows to all rows}
The unmarked density exponent supplied by the existing source is
\begin{equation}\label{eq:rho}
 \rho(\delta)=1-\frac{10\delta}{3(5-2\delta)},
 \qquad \frac7{10}\le\delta\le\frac34.
\end{equation}
The adjective ``unmarked'' means that this count requires no positive
prime-slot supply. It does not mean that arbitrary moving ray characters
can be inserted. A single fixed finite character list is selected before
the live labels.

Every nonzero row has a sixth-power decomposition $j=kb^6$, with the
source's primary-generator and unit conventions. The extra factors of
$b$ leave Euler holes in the displayed character. In the positive
half-plane they preserve the zeros of its primitive inducing character.
This permits the following general counting step.

\begin{lemma}[Summing the sixth-power part]\label{lem:sixth}
Assume a uniform count $\ll Y^{\rho+\eps}$ for the relevant
sixth-power-free cores of norm at most $Y$, for every $Y\ge1$.
If $\rho\ge2/7$, the corresponding count of all rows of norm at most
$H\ge1$ is $\ll H^{\rho+\eps}$. The constant may be chosen uniformly
for $\rho$ in a fixed interval with lower endpoint $2/7$, provided the
core estimate is itself uniform.
\end{lemma}
\begin{proof}
Up to the fixed unit-sector convention, the count is at most
\[
 H^{\rho+\eps}\sum_{Nb\le H^{1/6}}(Nb)^{-6(\rho+\eps)}.
\]
The ideal series converges since $6\rho\ge12/7>1$, and a single
convergent series with exponent $12/7$ dominates the sum uniformly.
If the input estimate was initially only eventual, bounded windows must
first be paid with a uniform constant. This is a genuine requirement,
not an automatic consequence of an eventual estimate with varying data.
\end{proof}

Here $\rho(7/10)=19/54$ and $\rho(3/4)=2/7$, and $\rho$ is decreasing.
Fixed-inducing presentations stay excluded from this count; their
exceptional contribution must be treated separately.

\subsection{A simple branch that explains the saving}
Consider the second-transform branch in which the same squarefree label
$P$ occurs once on each side, with $\log_Z NP=q$. Its actual character
row is
\[
 j=P^2h',\qquad \log_Z Nj\le M-6a+o(1).
\]
At a fixed $j$, the labels $P$ satisfy $P^2\mid j$. There are only
divisor-many such labels, at an arbitrarily small power cost. Thus a
density bound on $j$ controls the joint population $(P,h')$ without
an additional factor $Z^q$.

This ordering matters. If one first freezes $P$, bounds the remaining
rows for each twist, and finally pays for all choices of $P$, one loses
the information that $P$ is already contained in $j$.

For this branch the source normalization leads to the proposed exponent
\begin{equation}\label{eq:simplehigh}
 A-M+5a+(M-6a)\rho(\delta)+\delta(A-q).
\end{equation}
Relative to the allowance $A$, it has excess
$M(\rho-1)+\delta A+a(5-6\rho)-\delta q$.
Equation~\eqref{eq:simplehigh} illustrates the analytic accounting in the
source note; the full transformed identity and all masks remain part of
the analytic obligation. The illustration alone does not prove
Hypothesis~\ref{hyp:local}.

\subsection{The effective integer row in the general branch}
The first common-support extraction introduces widths $c,d,p$, where
$p$ is now the first common radical width. Let $R$ be its active
conductor width and $E$ the mask-divisor width. These are local
bookkeeping variables; the $p$ of Section~\ref{sec:endpoint} will instead
denote a detector parameter. Set $T=(c-d)_+$.

The signs of the two separated twists are opposite. We fix the convention
of the formal source: $t_{\mathfrak p}$ is the left multiplicity minus
the right multiplicity, modulo six. The first coefficient contains the
active character on the left and its conjugate on the right. After
reciprocity, the candidate row residues are therefore
\[
 e_C(\mathfrak p)=(1+t_{\mathfrak p})\bmod6,\qquad
 e_D(\mathfrak p)=(1-t_{\mathfrak p})\bmod6.
\]
Draft 1 used the opposite sign convention without aligning it explicitly
with the formal source. The scalar bounds below are unchanged, since
they treat either residue in the same way; the source character identity
must nonetheless respect this choice.
Define
\[
 F_C=\prod_{\mathfrak p\mid r}\mathfrak p^{e_C(\mathfrak p)},\quad
 L_C=\log_Z NF_C,\quad
 R_{\rm zero}=\sum_{\substack{\mathfrak p\mid r\\e_C(\mathfrak p)=0}}
                       \log_Z N\mathfrak p.
\]
An exponent zero here records cancellation of a phase; its old puncture
is still retained.

The proposed reindexing after the second common-support extraction is
\begin{equation}\label{eq:integerrow}
 w=F_CeG_cV_{\rm id}h',\qquad
 \log_Z Nw\le M-R-T+L_C-6a+o(1).
\end{equation}
The factors $e,G_c,V_{\rm id}$ are the source's mask and second
extraction factors. The analytic row ball must actually be enlarged to
the smooth envelope
\begin{equation}\label{eq:envelope}
 K_0+(d-c)_++6a,\qquad K_0=2A-c-d+R+E-M,
\end{equation}
before the second transformation. One cannot obtain the smaller dual
range in \eqref{eq:integerrow} merely by declaring a larger bound for an
unchanged first row ball.

At fixed $w$, the intended reconstruction makes the active first primes,
the mask primes, the second common primes, and the first residual
projector label divisor-many. First common primes absent from $w$ must
still be counted; their width is
\begin{equation}\label{eq:missing}
 p-R-E+R_{\rm zero}.
\end{equation}
In particular, a canceled prime cannot simultaneously be omitted from
$w$ and treated as one of its divisors. A general arithmetic justification
of this counting step is now given in the companion reconstruction note
\cite{gaov2}. At a fixed row $w$, write $L$ for the radical of first
common primes absent from $w$. All the remaining labels are supported
on $\operatorname{rad}(w)L$. For a fixed bound on their logarithmic norms,
Rankin's inequality gives, for every $\eps>0$,
\[
 \#\{I:NI\le Z^T,\ \operatorname{rad}I\mid Q\}
 \ll_{B,T,\eps}Z^\eps\qquad(NQ\le Z^B).
\]
Indeed the count is at most
$Z^{aT}\prod_{\mathfrak p\mid Q}(1-(N\mathfrak p)^{-a})^{-1}$;
take $aT\le\eps/2$ and bound the product by a constant times
$(NQ)^{\eps/(2B)}$. The finite exceptional set of small primes contributes
only a constant. Counting the missing radical costs its norm exponent
\eqref{eq:missing}. This proves a general reconstruction bound under
the stated factorization and allocation assumptions, not just for a
finite sample of valuations. Proving those assumptions for every actual
source allocation remains open.

The first labels and the full coprimality projector remain inside the
positive norm until this joint count is applied. Their M\"obius signs
are retained through the exact allocation. A common Fourier measure is
chosen on the ambient parameter box before labels and bins; only its
evaluation points depend on the later arithmetic data. This order is
needed to keep independent heights and smooth-profile bounds uniform.

\subsection{What the proposed accounting gives}
Let $b_2,g_2,t_2$ denote the second common widths, with
$0\le g_2\le b_2$ and $t_2\ge0$, and put
$B_c=(3c-5d-R)_+/6$. Relative to the joint allowance
$p+A-c+R+E+B_c$, the proposed excess is
\begin{align}\label{eq:ledger}
 \mathcal E={}&M(\rho-1)+\delta A+a(5-6\rho)\notag\\
 &+(1-\rho)T-\delta c-(1+\rho)R-E-B_c
      +\rho L_C+R_{\rm zero}\notag\\
 &+g_2-t_2-(1+\delta)b_2+\eta.
\end{align}
The rational audit checks this algebraic identity from the encoded
source exponents. It does not prove that the analytic expression being
estimated is equal to those encoded exponents.

The favorable signs can be checked directly. Using
$c=\min(c,d)+T$, $\min(c,d)\ge p$, $R\le p$, and
$L_C\le5(R-R_{\rm zero})$, the second line of \eqref{eq:ledger} is at
most
\[
 (4\rho-1-\delta)p-E-B_c
 \le-\frac{79}{270}p-E-B_c.
\]
Indeed $\rho+\delta-1>0$, $5\rho-1>0$, and $4\rho-1>0$ throughout
the interval. Also $4\rho-1-\delta$ is decreasing and equals
$-79/270$ at $\delta=7/10$. The third line without $\eta$ is at most
$-\delta b_2-t_2\le0$.

Finally \eqref{eq:window} implies the weaker bound
$A/M\le1-(16/5)(3/100)=113/125$. Since
\[
 \rho'(\delta)=-\frac{50}{3(5-2\delta)^2},
\]
the function $\rho(\delta)+(113/125)\delta$ decreases on
$[7/10,3/4]$. Therefore
\begin{equation}\label{eq:highreserve}
 1-\rho(\delta)-\delta A/M\ge\frac{259}{16875},\qquad
 5-6\rho(\delta)\le\frac{23}{7}.
\end{equation}
For $a=M/2000$ and $\eta\le M/1000$, equations
\eqref{eq:ledger}--\eqref{eq:highreserve} give
\[
 \mathcal E\le
 -\frac{259M}{16875}+\frac{23M}{14000}+\frac M{1000}
 <-\frac M{100}.
\]
This is a substantial scalar reserve. Its usefulness depends on proving
the analytic reconstruction and normalization preceding it.

\section{Other branches and uniformity}\label{sec:interfaces}

\subsection{Comparison terms and additional masks}
The centered source uses equal-product comparison rectangles with one
plain length $M/4$. For an additional common mask, the proposed argument
expands that mask only while bounding the comparison terms. After
reflecting the second natural daughter if its length exceeds $M/2$, the
two plain lengths sum to at most $3M/4+\xi$. Their existing-moment
budget is consequently at most
\[
 \frac{3M}{4}+\frac92\frac M{25}+\xi
 =\frac{93M}{100}+\xi<M
\]
for sufficiently small padding $\xi$. This pays comparison daughters
without asserting that they belong to the local window
\eqref{eq:window}. The mask remains on both rectangles of the original
centered expression.

\subsection{The amplifier errors}
For the three extracted valuation alternatives $1,6,7$, the relevant
width pairs are $(\lambda,\lambda)$, $(6\lambda,0)$, and
$(7\lambda,\lambda)$, with $\lambda\ge\sigma/4$. The old moment
pays these branches. The source's defect estimate is
\[
 F_{\rm act}\le g+\sigma+w-w_o+\eta,
\]
its base cost is at most $F_{\rm act}/6$, and its saving is at least
$w$. Substituting $g\le M/10000$ gives the following lower margins
before padding:
\begin{center}
\begin{tabular}{ccc}
\toprule
Valuation & $(w,w_o)$ & Saving minus base cost\\
\midrule
$1$ & $(\lambda,\lambda)$ & $13M/120000$\\
$6$ & $(6\lambda,0)$ & $193M/120000$\\
$7$ & $(7\lambda,\lambda)$ & $119M/60000$\\
\bottomrule
\end{tabular}
\end{center}
The normalized prime average must be retained. Introducing a second
count of the auxiliary primes would change these margins.

\subsection{Exceptional terms cannot be hidden in an error}
Lemma~\ref{lem:parent} gives $A-M<-M/10$. The source's second
zero-frequency term has excess at most $A-M+5\sigma/3+\eta$, and
the centered fixed-inducing calculation retains a maximum $A-M$ before
its terminal loss $3\sigma$. Both fit within this negative margin.
Replacing $A-M$ by $(A-M)_+$ would discard the saving that pays the
fixed amplifier.

For the enlarged row ball \eqref{eq:envelope}, the first exceptional
reserve must also be checked again. The proposed expressions are
\[
 F_1=\begin{cases}
 (c+5d+R+E)/6+B_c-5a,&c\ge d,\\
 c+(R+E)/6+B_c-5a,&c\le d.
 \end{cases}
\]
They imply $F_1\ge2c/3-5a$, while the second reserve remains
$F_2\ge2b_2/3$. The needed scalar inequalities hold; the exact
centered cancellation, forced congruences, and raw extraction factors
must still be carried through the physical proof.

Uniformity is part of the assertion. One first fixes the requested
final loss, then chooses the aggregate small losses below both that
loss and the margin in \eqref{eq:lowreserve}. Profile orders, height
orders, and the eventual threshold are chosen through the finite chain
of transformations. A loss per slot, or a character family chosen only
after the live labels, would not supply the required estimate.

\section{The exact conditional endpoint}\label{sec:endpoint}

This section describes an independently reproducible scalar conclusion.
It uses the existing detector function; only the physical moment
coefficient is changed. Here $p\in[0,1/2]$ is a detector parameter,
unrelated to the common radical width in Section~\ref{sec:high}.

\subsection{The rational geometry and the exponent function}
Set
\begin{align}\label{eq:geometry}
 \ell&=\frac{626701}{3750000}, &
 h&=\frac{3969737987}{5000000000},\notag\\
 x&=1+\ell-h=\frac{5597590039}{15000000000}, &
 y&=h-2\ell=\frac{6895605961}{15000000000},\notag\\
 \kappa_*&=2\theta-1=\frac{1874433}{2500000}, &
 C&=\frac{2249}{500}.
\end{align}
Then $\theta=11/12-\ell/4$ and $x+y=1-\ell$.
For $\kappa\ge\kappa_*$ define
\begin{align}\label{eq:detector}
 D_\kappa(p)&=3-p-\frac{2p}{3\kappa}, &
 B_\kappa(p)&=2(1-p)\left(1-\frac{p}{3\kappa}\right),\notag\\
 c_\kappa(p)&=B_\kappa(p)/D_\kappa(p), &
 V&=\frac56-\delta,\notag\\
 R_\kappa(\delta,p)&=1-\delta+
       \frac{\delta c_\kappa(p)V}{2(V+\delta c_\kappa(p))}.
\end{align}
We use a lower-case $c_\kappa$ to distinguish the detector coefficient
from the physical coefficient $C$. The candidate upper-row exponent is
\begin{equation}\label{eq:F}
 F_{C,\kappa}(\delta,p)=\frac{1+\ell}{2}+\frac\delta2
 +h\left(\frac{3R_\kappa(\delta,p)}4-\frac5{12}\right)
 +\delta p\left(\ell-\frac{h-4\ell}{2C}\right).
\end{equation}

\begin{proposition}[Exact scalar certificate]\label{prop:certificate}
With \eqref{eq:geometry},
\[
 F_{C,\kappa_*}(\delta,p)\le\theta-10^{-7}
 \quad\text{for every }(\delta,p)\in[0,3/4]\times[0,1/2].
\]
Moreover $\theta<b_0$ for the exact boundary \eqref{eq:quartic}.
\end{proposition}
\begin{proof}[Computer-assisted proof of the first assertion]
All constants in this calculation are rational. Put
$J=VD_{\kappa_*}+\delta B_{\kappa_*}$. Then
\[
 R_{\kappa_*}=\frac{(1-\delta)J+delta B_{\kappa_*}V/2}{J}.
\]
The stored certificate first gives $J\ge1/10$ on the whole rectangle.
After multiplying $\theta-10^{-7}-F_{C,\kappa_*}$ by $J$, it suffices
to verify a polynomial inequality. The audit partitions the rectangle
into 18 rational subrectangles and writes the polynomial in the
Bernstein basis on each one. Every coefficient is nonnegative; the
recorded closed subrectangles cover the entire original domain.
Appendix~\ref{app:reproduce} identifies the executable audit and its
exact certificate. This is an exact finite verification, not a sampling
of the exponent on a numerical grid.
\end{proof}
\begin{proof}[Proof of the comparison with the quartic]
We have $P(1/6)<0$, and exact substitution gives
\[
 P(\ell)=\frac{3484615068030068970503}
                    {21972656250000000000000000}>0.
\]
For $X\in[1/6,1/5]$,
\[
 P'(X)=3708X^3-9405X^2+4866X+275
 \ge-\frac{9405}{25}+\frac{4866}{6}+275
 =\frac{3549}{5}>0.
\]
Thus $\ell>\ell_0$, and $\theta=11/12-\ell/4<b_0$.
\end{proof}

\subsection{Why the local length restriction matters}
The row exponent in the numerator varies; it is not always $h$.
Let $d$ denote that exponent and $t$ a normalized free Euler shift.
The local estimate is intended only for
\begin{equation}\label{eq:applicationwindow}
 h-10^{-3}\le d\le h+10^{-10},\qquad 0\le t\le10^{-3}.
\end{equation}
The audit checks that both reflected plain lengths, normalized by the
row budget, lie strictly between $0.41$ and $0.44$, and that the selected
prime width lies strictly between $0.03$ and $0.04$. It also checks
that the old $9/2$ defect is less than $1/12500$ and that the original
physical slots supply the selected product.

For clarity, a general scalar row exponent used by the audit is
\begin{align}\label{eq:G}
 G_{c,\kappa}(d;\delta,p)={}&\frac12
 +h\left(\frac{17}{50}-\frac16\right)-\frac y2-\frac\ell2
 +d\left(\frac14-\frac{17}{50}\right)+\frac\delta2\notag\\
 &+\frac{3d}{4}R_\kappa(\delta,p)
 +\delta p\left(\ell-\frac{2y-d}{2c}\right).
\end{align}
It is affine in $d$ on the range in question, since $d\ge1/2>y$.
At $d=h$ and $c=C$ it reduces to \eqref{eq:F}. The exact checks are
\begin{center}
\begin{tabular}{lccc}
\toprule
Row endpoint & Coefficient $c$ & Saving below $\theta$ & Rectangles\\
\midrule
$h$ & $C$ & $10^{-7}$ & 18\\
$h-10^{-3}$ & $6\kappa_*$ & $10^{-5}$ & 6\\
$1/2$ & $6\kappa_*$ & $10^{-3}$ & 1\\
\bottomrule
\end{tabular}
\end{center}
The lower interval therefore uses the old estimate. At the junction,
$C<6\kappa_*$ and $2y-d>0$, so the local exponent is no larger
than the old one. Affinity bounds the intermediate local rows. Above
$h$, a scalar slope bound of $2$ leaves saving
$10^{-7}-2\cdot10^{-10}>9\cdot10^{-8}$.

For larger free shifts the source's existing negative coefficient
contribution must be kept. Write $q=\delta p$, so $0\le q\le\delta/2$.
The audit verifies, for the normalized available width in this window,
\[
 -\frac{\delta t}{2}-\frac{qz_{\rm old}(t)}2
 \le-\frac{qz_{\rm new}(0)}2,\qquad t\ge10^{-3}.
\]
Here $z_{\rm old}(t)=\max(0,b-2t)/c$ with
$6\kappa_*\le c\le9/2$, and $z_{\rm new}(0)=b/C$; $b$ is the
residual budget after the unshifted reflected plain factors. The
normalized reserve at the first shift exceeds $3/10000$ and increases
thereafter, including after the old available width becomes zero.
These are scalar interface checks. They do not by themselves show that
every physical Euler term has been represented by this interface.

\subsection{The intended noncircular conclusion}
Suppose Hypothesis~\ref{hyp:local} is proved and can be inserted into
the original weighted numerator with all coefficients, masks, conductor
losses, principal terms, and family reductions preserved. One must then
run the existing zero-supremum argument for that assembled numerator.
Let $\beta_*$ be the supremum used in the source, with its sentinel
$1/2$. The $7/8$ initialization gives $\beta_*\le7/8$. If
$\beta_*>\theta$, put $\kappa=2\beta_*-1$; then
$\kappa_*<\kappa\le3/4$.

For the new upper-row expression, only the detector depends on $\kappa$.
The source bounds $\partial_\kappa c_\kappa<1/18$ and
$\partial_cR\le5/12$ give
\[
 \partial_\kappa F_{C,\kappa}
 \le\frac{3h}{4}\frac5{12}\frac1{18}
 =\frac{3969737987}{288000000000}<\frac1{50}.
\]
Consequently Proposition~\ref{prop:certificate} implies
\begin{equation}\label{eq:feedback}
 F_{C,2\beta_*-1}(\delta,p)-\beta_*
 \le-10^{-7}-\frac{24}{25}(\beta_*-\theta)<0.
\end{equation}
The old rows retain their existing feedback estimates. Once all analytic
pieces and the small losses are supplied, such strict savings are the
input to the source's contradiction to the definition of $\beta_*$.
The hypothesis on the slot zero bound is an equality at the chosen
$\kappa$; the desired improved zero-free conclusion is not used as an
input. Equation~\eqref{eq:feedback} explains the intended mechanism,
not a completed proof that $\beta_*\le\theta$.

\section{What has been verified, and what remains}\label{sec:gaps}

The submitted package contains exact source calculations, exact endpoint
certificates, and nine Lean components. The source audit checks 11
continuous rational sign statements, symbolic exponent identities,
317,376 finite joint integer-fiber cases, and 72 sextic CRT cases,
including negative controls. Finite cases are useful checks on a
proposed reconstruction; they do not prove it for all arithmetic data.

The component roles are summarized below. Words such as ``implies'' are
important: several results still have input estimates or witness
conditions as explicit hypotheses.
\begin{center}\small
\begin{tabularx}{\textwidth}{@{}l>{\raggedright\arraybackslash}X@{}}
\toprule
Component & Proved role within the component\\
\midrule
\code{FiniteDeletionBins} & Finite Euler deletion preserves the relevant positive-half-plane zeros and buffered maxima.\\
\code{UnmarkedDensity} & Unmarked inverse and plain moment inputs imply the fiber density exponent with explicit losses.\\
\code{BufferedSlot} & The natural relative slot satisfies a buffered pointwise bound with punctures and independent frequencies.\\
\code{RestrictedIncumbent} & Slot removal on selected low rows feeds the existing remaining-state moment, with separate costs.\\
\code{SixthRowMask} & The sixth-power row decomposition preserves the original Euler holes and zero-bin data.\\
\code{RawSixthDensity} & An all-window core density estimate implies the raw-row count by ideal convolution.\\
\code{UniformDensity} & Uniform eventual estimates extend to bounded windows, and the convolution constant can be uniform in the exponent.\\
\code{UnmarkedBatch} & Actual witness labels and two dyadic lengths partition the unmarked rows; supported witnesses suffice.\\
\code{ActualUnmarkedInputs} & Supplies the incumbent inverse profiles for a fixed finite family, constructs witnesses from actual buffered maxima, and pays the fixed-label and dyadic costs.\\
\bottomrule
\end{tabularx}
\end{center}
The second replay compiled all nine components with Lean 4.34.1 against
cached dependencies. The axiom audit covered sixteen named declarations.
Its reported axioms were \code{propext}, \code{Classical.choice}, and
\code{Quot.sound}. Imported dependencies were not rebuilt, and canonical
target comparison and independent three-kernel checks remain pending.

The outstanding mathematical work is concrete:
\begin{enumerate}
\item Supply the remaining plain unmarked moment through the actual
zero-slot source and eligibility adapter. The inverse profile estimates,
the fixed-label and dyadic costs, and a constructor from actual buffered
maxima are now supplied. The application must still prove that the real
row family meets the constructor's nonprincipal, modulus, coefficient,
and zero-bin conditions, uniformly in conductor and height.
\item Prove the signed-twist identification \eqref{eq:integerrow} and
the joint reconstruction leading to \eqref{eq:missing}, for both sides
and every projector allocation, retaining all natural zeros.
\item Complete the comparison-only mask reflection and exceptional
terminal estimates at the actual enlarged row radius. Keep the negative
margins that pay the fixed amplifier.
\item Assemble \eqref{eq:moment} with the common Fourier measure,
independent heights, all required profiles, and the original physical
state. Then carry the full numerator through powerful-conductor
deficits, free Euler shifts, principal signals, and the family reductions.
\item Prove the exact ordinary Dirichlet, zeta, and Eisenstein-Hecke
targets and obtain independent canonical comparison and kernel
verification. Component compilation does not supply these targets.
\end{enumerate}
There is a specific interface obstruction at the first item. The
upstream \path{source_batch_plain_unmarked} theorem takes the original
marked batch, its large-row restriction, and its initial-state admission.
An empty selected product does not remove those hypotheses. A proof for
the present unmarked family must generalize that source adapter and pay
small-core rows uniformly. The present submission therefore remains a
research candidate with missing mathematical assembly, not merely a
missing verification receipt.

\section{Attribution and source relationship}\label{sec:credits}

OpenAI's development supplies the initial $7/8$ theorem and the analytic
and arithmetic framework. Levy's work supplies the weighted-numerator
input adapted by the Cycle25 source. Collet's work supplies the variable
moment, detector, quartic optimization, and source architecture used
throughout this draft. Our proposed contribution is the local
$4.498$ refinement, its joint-label analysis, and the conditional exact
endpoint calculation. These descriptions do not imply endorsement by
any upstream author or maintainer.

This manuscript is an exposition of the three candidate notes in the
research package \cite{gao} and the subsequent input and reconstruction
work \cite{gaov2}, with definitions and limitations made
explicit. It does not change their verification status. Research,
code, formalization, and writing received AI assistance; the internal
reviews recorded with the package are model-assisted development
reviews, not independent human referee reports. The package retains its
Apache-2.0 license and upstream notices.

\appendix
\section{Reproduction and the certificate format}\label{app:reproduce}

The package linked in \cite{gao} contains three notes under the directory
\path{research/framework/programs/fixed_strip/candidates/}:
\begin{itemize}
\item \code{LOCAL\_PHYSICAL\_MOMENT\_EXTENSION.md};
\item \code{JOINT\_HIGH\_ZERO\_SQUARED\_LABEL\_PAYMENT.md};
\item \code{LOCAL\_MOMENT\_TRACKER\_REFINEMENT.md}.
\end{itemize}
The first explains the proposed local moment, the second isolates a
useful squared-label branch, and the third records the endpoint
geometry. In particular, the coefficient $21/5$ used illustratively in
the second note is not the coefficient claimed in
Hypothesis~\ref{hyp:local}, and is not an established stronger moment.

From the submitted package root, with Python and the recorded SymPy
version available, run
\begin{verbatim}
python3 reproduce.py --fetch-pinned-paper
python3 research/tools/local_physical_moment_audit.py --check
python3 research/tools/local_moment_endpoint_audit.py --check
\end{verbatim}
The fetched manuscript is a pinned input, checked by its stored digest;
it is not a paper authored by this project. The outputs in
\code{research/numerics/results/} record the exact assertions and their
scope. Both reports explicitly deny verification of the new full
analytic theorem. The component replay has separate compiler and
dependency prerequisites documented in its own README.

For completeness, the Bernstein certificate used in
Proposition~\ref{prop:certificate} is based on the identity
\[
 f(u,v)=\sum_{i=0}^{m}\sum_{j=0}^{n}
 b_{ij}\binom mi u^i(1-u)^{m-i}\binom nj v^j(1-v)^{n-j},
 \qquad 0\le u,v\le1.
\]
The basis functions are nonnegative and their sum is one. Thus
$b_{ij}\ge0$ for all $i,j$ proves $f\ge0$ on the whole square.
An exact affine substitution maps each recorded rational rectangle to
that square. Rational reconstruction of the coefficients and exact
coverage checks make the conclusion inspectable without trusting
floating-point optimization.

\begin{thebibliography}{9}
\raggedright
\bibitem{collet}
Hailey Collet, \emph{The quartic QRH boundary}, 10 October 2026.
Cycle25 manuscript and source, revision
\code{58344dfdbe756cf2f743da1908ffe0582418cf5b}.
\href{https://github.com/HaileyStorm/QuasiRiemannTracker/blob/58344dfdbe756cf2f743da1908ffe0582418cf5b/public/proofs/cycle25-quartic-20261010/paper.tex}{Pinned manuscript}.
Equation references B0 and B3 in the companion notes refer to this revision.

\bibitem{openai}
OpenAI, \emph{OpenAI/math}, formal mathematical development,
revision \code{fd4aeeb2ee4fc729c18d98444fed42fd0529eeeb}.
\href{https://github.com/openai/math/tree/fd4aeeb2ee4fc729c18d98444fed42fd0529eeeb}{Pinned source repository}.

\bibitem{levy}
Akash Levy, \emph{Weighted-numerator proof}, source package,
revision \code{2fd60c0926b66ea18d7436f5ed55250fd006ab5d}.
\href{https://github.com/akashlevy/QuasiRiemannTracker/tree/2fd60c0926b66ea18d7436f5ed55250fd006ab5d/proofs/akashlevy-20261009-weighted-numerator}{Pinned proof package}.

\bibitem{gao}
Zhongpai Gao, with AI assistance,
\emph{Local physical moment research candidate: 0.8748866},
11 October 2026. Source-review snapshot,
revision \code{7f9166e2d626586cc3c89853497c0e0e22e856ee}.
\href{https://github.com/Gaozhongpai/QuasiRiemannTracker/tree/7f9166e2d626586cc3c89853497c0e0e22e856ee/proofs/gao-local-moment-20261011}{Research notes, certificates, and Lean components};
\href{https://github.com/GTimG/QuasiRiemannTracker/pull/15}{draft pull request 15}.

\bibitem{gaov2}
Zhongpai Gao, with AI assistance,
\emph{Joint-label reconstruction and actual unmarked inputs},
11 October 2026, second research revision.
\href{https://github.com/Gaozhongpai/QuasiRiemannTracker/tree/gao-local-moment-source-review-20261011/proofs/gao-local-moment-20261011-v2}{Companion note and nine compiled components}.
The general counting lemma retains the physical factorization as an
explicit application obligation.
\end{thebibliography}
\end{document}
